\section{Transformation formulas}\label{sec:transformations}

\subsection{A finite pullback--gauge certificate}

\begin{lemma}[Pullback certificate]\label{lem:pullback-certificate}
Let \(v(w)\) satisfy \(\mathcal L_{A,B,C}v=0\).  Let
\(w=\phi(z)\) and \(g(z)=p(z)v(\phi(z))\), where \(p,\phi\) are rational or
algebraic named functions.  After differentiating twice and replacing
\(v''\) with the expression supplied by its differential equation, one gets
\[
 \mathcal L_{a,b,c}g=R_0(z)v(\phi(z))+R_1(z)v'(\phi(z)).
\]
If exact algebraic certificates prove \(R_0=R_1=0\), and if
\(g(0)=1\), then
\[
 g(z)=\F a b c z
\]
as a formal series.  On any branch domain with effective convergence, the
identity holds between named analytic functions.
\end{lemma}

\begin{proof}
The chain rule is a finite symbolic calculation.  Clearing all denominators
turns \(R_0=R_1=0\) into polynomial identities in the parameter and pullback
codes.  The resulting \(g\) satisfies the target ODE.  Its constant term fixes
the normalized regular solution by \cref{cor:ode-unique}.  If a resonant
parameter makes the regular recurrence singular, the same certificate is
read after localization or at each finite truncation where the required
denominators are guarded.
\end{proof}

\begin{remark}
This lemma is a practical proof checker for transformation formulas.  It does
not rely on recognizing two analytic functions after the fact.  The
certificate consists only of two cleared rational-function numerators and a
normalization value.
\end{remark}

\subsection{Euler and Pfaff transformations}

\begin{theorem}[Euler transformation]\label{thm:euler-transform}
On the normalized branch at \(z=0\),
\begin{equation}\label{eq:euler-transform}
 \F a b c z=(1-z)^{c-a-b}\F{c-a}{c-b}{c}{z}.
\end{equation}
\end{theorem}

\begin{proof}
Apply \cref{lem:pullback-certificate} with \(\phi(z)=z\),
\(p(z)=(1-z)^{c-a-b}\), and source parameters
\((c-a,c-b,c)\).  Using
\(p'/p=-(c-a-b)/(1-z)\), direct collection makes both residual
coefficients zero.  The prefactor and inner series have constant term \(1\).
The calculation is polynomial after multiplication by \((1-z)^2\).
\end{proof}

\begin{theorem}[Pfaff transformations]\label{thm:pfaff}
On the normalized branch at zero,
\begin{align}
 \F a b c z
 &=(1-z)^{-a}
   \F a{c-b}c{\frac{z}{z-1}},
   \label{eq:pfaff-a}\\
 \F a b c z
 &=(1-z)^{-b}
   \F{c-a}b c{\frac{z}{z-1}}.
   \label{eq:pfaff-b}
\end{align}
\end{theorem}

\begin{proof}
For \eqref{eq:pfaff-a}, use
\(\phi=z/(z-1)\), \(p=(1-z)^{-a}\), and parameters
\((a,c-b,c)\) in \cref{lem:pullback-certificate}.  The derivative data
\[
 \phi'=-\frac1{(z-1)^2},\qquad
 \phi''=\frac2{(z-1)^3},\qquad
 \frac{p'}p=\frac a{1-z}
\]
make the residual numerators vanish.  The second formula follows by symmetry
in \(a,b\).  All denominators are powers of \(1-z\), so the formal identity
is checked in a localized polynomial ring and the analytic identity needs a
branch disk avoiding \(1\).
\end{proof}

\begin{corollary}[Half-argument conversion]\label{cor:half-minus-one}
\[
 \F a b c{1/2}=2^a\F a{c-b}c{-1}
\]
whenever both named values and the chosen power of two are defined.
\end{corollary}

\begin{proof}
Put \(z=1/2\) in \eqref{eq:pfaff-a}.
\end{proof}

\subsection{Quadratic transformations}

\begin{theorem}[A standard quadratic transformation]\label{thm:quadratic}
On the branch normalized at zero,
\begin{equation}\label{eq:quadratic}
 \F a b{2b}{z}
 =(1-z)^{-a/2}
 \F{a/2}{b-a/2}{b+1/2}{\frac{z^2}{4z-4}}.
\end{equation}
\end{theorem}

\begin{proof}
Apply \cref{lem:pullback-certificate} with
\[
 \phi(z)=\frac{z^2}{4z-4},
 \qquad p(z)=(1-z)^{-a/2}.
\]
Use a source equation with parameters
\((a/2,b-a/2,b+1/2)\), and target parameters \((a,b,2b)\).
Multiplication by \(16(1-z)^4\) clears every derivative denominator.  The
coefficient of \(v'\) and then that of \(v\) expand to zero.  At zero,
\(p(0)=1\), \(\phi(0)=0\).  Moreover, the right derivative is
\(a/2=ab/(2b)\), so the normalization is compatible even in a formulation
using the first two coefficients.  This provides a finite checker rather
than an appeal to a table of transformations.
\end{proof}

\begin{theorem}[Kummer quadratic pair]\label{thm:kummer-quadratic}
For the normalized branches,
\begin{align}
 \F a b{1+a-b}{z}
 &=(1-z)^{-a}
 \F{a/2}{(a+1)/2-b}{1+a-b}
   {-\frac{4z}{(1-z)^2}},
 \label{eq:kummer-q1}\\
 \F a b{1+a-b}{z}
 &=(1+z)^{-a}
 \F{a/2}{(a+1)/2}{1+a-b}
   {\frac{4z}{(1+z)^2}}.
 \label{eq:kummer-q2}
\end{align}
\end{theorem}

\begin{proof}
Both are applications of \cref{lem:pullback-certificate}.  For the first use
\(\phi=-4z/(1-z)^2\), \(p=(1-z)^{-a}\); for the second use
\(\phi=4z/(1+z)^2\), \(p=(1+z)^{-a}\).  Clearing \((1\mp z)^4\) reduces the
claims to polynomial identities.  The constant term is \(1\).  The source
parameter relation \(c=1+a-b\) is exactly what cancels the residual
coefficient.
\end{proof}

\begin{remark}[Constructive domain statement]
The formal transformations do not need inequalities.  Their analytic reading
requires that the chosen disk avoid the poles of the pullback and that the
prefactor branch be named there.  The equality may then be continued along a
coded path by \cref{thm:path-continuation}.  This is more precise than a
single informal phrase such as ``by analytic continuation.''
\end{remark}

\subsection{The six local expressions and Kummer's 24 forms}

Let the six base expressions be
\[
 u_{0,1},u_{0,2},u_{1,1},u_{1,2},u_{\infty,a},u_{\infty,b}
\]
from \cref{prop:zero-solutions,prop:one-solutions,prop:infty-solutions}.
For any hypergeometric factor in one of these expressions, apply one of four
operations:
\begin{enumerate}[label=(\arabic*)]
\item leave it unchanged;
\item apply Euler's transformation;
\item apply the first Pfaff transformation;
\item apply the second Pfaff transformation.
\end{enumerate}
After simplifying the outer monomial, this produces the 24 forms tabulated by
Kummer and displayed by MathWorld.

\begin{theorem}[Kummer-form generator]\label{thm:kummer-24}
For generic nonresonant parameters, every expression generated by the
preceding \(6\times4\) construction solves the same hypergeometric differential
equation.  Within the local branch domain of its parent base expression, it
represents that same local solution.  The construction and its verification
are primitive recursive symbolic procedures.
\end{theorem}

\begin{proof}
Each base expression solves the target equation by the operator calculations
of \cref{sec:ode}.  Euler and Pfaff transformations replace only the inner
hypergeometric factor by a formally equal expression, so the parent function
is unchanged in its normalized local branch.  Syntactic simplification of
powers adds exponents and composes rational pullbacks; this is finite rational
arithmetic.  Alternatively, each of the 24 outputs can be checked directly
by \cref{lem:pullback-certificate}.  Resonant coincidences reduce the number
of distinct forms but do not invalidate the equation-solving claim.
\end{proof}

\begin{audit}[What ``all 24'' means]
The 24 entries are not 24 globally single-valued functions.  They are local
expressions containing powers such as \((-z)^{-a}\) and
\((1-z)^{c-a-b}\).  Equality of two entries away from the normalized local
disk includes a connection constant determined by continuation.  The formal
certificate proves the differential equation and local normalization; the
analytic branch statement carries a path and logarithm certificate.
\end{audit}

\subsection{Higher rational pullbacks}

\begin{proposition}[Finite checkability of algebraic transformations]
Suppose a proposed quadratic, cubic, or higher transformation has the form
\[
 \F a b c z=p(z)\F A B C{\phi(z)},
\]
where \(p\) and \(\phi\) are algebraic functions given by isolated polynomial
relations and local initial values.  If implicit differentiation supplies
named \(p',p'',\phi',\phi''\), then the identity has a finite
pullback--gauge certificate as in \cref{lem:pullback-certificate}.
\end{proposition}

\begin{proof}
Use the polynomial relations to eliminate all algebraic powers in the two
residual coefficients.  A Gröbner-style reduction is not foundationally
needed: a certificate explicitly expresses each cleared residual numerator as
a polynomial combination of the defining relations.  Checking that
combination is primitive recursive polynomial arithmetic.  The initial value
selects the intended local branch.
\end{proof}

This proposition covers the specific transformations stated on the source
pages and gives a direct formalization route for the many additional Goursat
transformations that those pages mention but do not print.
